% !TeX root = ../navier-stokes-manuscript.tex
% ============================================================
% 2.6 Finite Adjacent Rectangles
% ============================================================

\subsection{Finite Adjacent Rectangles}\label{sub:FiniteAdjacentRectangles}

Only finite estimates can build the intersection.  Begin on
\(I=(0,T)\), \(T<T_*\), with the first request that needs adjacent rows: bring
both
\[
  V_0=u
  \qquad\text{and}\qquad
  V_1=\partial_tu
\]
to one time slice in \(H^2(\R^3)\).  The upper spatial row
\(V_0\in L^2_tH^3_x\) alone has no time trace.  Its lower adjacent row is
\[
  \partial_tV_0=V_1\in L^2(I;H^1).
\]
To trace \(V_1\) at that same height, repeat the pair once:
\[
  V_1\in L^2(I;H^3),
  \qquad
  \partial_tV_1=V_2\in L^2(I;H^1).
\]
Together these are
\[
  V_0,V_1\in L^2(I;H^3),
  \qquad
  V_1,V_2\in L^2(I;H^1).
\]
The repeated \(V_1\) plays two roles: it closes the lower side of the first
trace and occupies the upper side of the second.

The reason the levels straddle \(H^2\) is visible directly.  Let
\[
  A:=(1-\Delta)^{1/2},
  \qquad
  \norm{v}_{H^s}=\norm{A^sv}_2.
\]
For smooth \(v\),
\[
  \norm{v(t)}_{H^2}^2-\norm{v(s)}_{H^2}^2
  =
  2\operatorname{Re}
  \int_s^t
  \left\langle
    A\partial_\tau v(\tau),
    A^3v(\tau)
  \right\rangle_{L^2}
  \,d\tau.
\]
The integrand is bounded by
\[
  \left|
    \left\langle
      A\partial_\tau v,
      A^3v
    \right\rangle_{L^2}
  \right|
  \le
  \norm{\partial_\tau v}_{H^1}
  \norm{v}_{H^3},
\]
which is integrable in time whenever
\[
  v\in L^2(I;H^3),
  \qquad
  \partial_tv\in L^2(I;H^1).
\]
Approximation and weak continuity consequently give
\[
  v\in C\bigl([0,T];H^2(\R^3)\bigr).
\]
Applying this once to each adjacent pair yields
\[
  V_0,V_1
  \in
  C\bigl([0,T];H^2(\R^3)\bigr).
\]
Collect the four norms consumed by these two traces in
\begin{align*}
  \mathfrak R_{1,2}(u;T)
  :={}&
  \norm{u}_{L^2(0,T;H^3)}^2
  +
  \norm{\partial_tu}_{L^2(0,T;H^3)}^2
  \\
  &+
  \norm{\partial_tu}_{L^2(0,T;H^1)}^2
  +
  \norm{\partial_t^2u}_{L^2(0,T;H^1)}^2.
\end{align*}
All four entries belong to the original solution.  Their grouping is a record
of what the two traces consume, not a construction of four new fields.

Now choose independent finite cutoffs \(N,M\in\Nzero\).  To trace each retained
temporal row at \(H^M\), use the adjacent scale
\[
  H^{M+1}(\R^3)
  \hookrightarrow
  H^M(\R^3)
  \hookrightarrow
  H^{M-1}(\R^3),
\]
with \(H^{-1}\) as the lower member when \(M=0\).  For every
\(0\le n\le N\), the required pair is
\[
  V_n\in L^2(0,T;H^{M+1}),
  \qquad
  V_{n+1}=\partial_tV_n\in L^2(0,T;H^{M-1}).
\]
Their squared finite velocity norm is
\begin{align*}
  \mathfrak R_{N,M}(u;T)
  :={}&
  \sum_{n=0}^{N}
  \norm{V_n}_{L^2(0,T;H^{M+1})}^{2}
  \\
  &+
  \sum_{n=0}^{N}
  \norm{V_{n+1}}_{L^2(0,T;H^{M-1})}^{2}.
\end{align*}
For a smooth representative, each row satisfies
\[
  \left|
    \frac{d}{dt}\norm{V_n(t)}_{H^M}^{2}
  \right|
  \le
  2
  \norm{V_{n+1}(t)}_{H^{M-1}}
  \norm{V_n(t)}_{H^{M+1}}.
\]
Its right-hand side is integrable, so density and weak continuity yield
\[
  V_n
  \in
  C\bigl([0,T];H^M(\R^3)\bigr),
  \qquad
  0\le n\le N.
\]
The temporal cutoff \(N\) chooses how many responses are retained; the spatial
cutoff \(M\) chooses their common trace height.  They remain independent, and
the rectangle constant may depend on either one.

The scalar \(\mathfrak R_{N,M}(u;T)\) remembers only the velocity norms.  A
rectangle must also remember why those rows belong together: retain the
normalized pressure coordinates and the differentiated equations.  For the
pair generated by \(u_0\), define
\[
  \cR_{N,M}[u_0;T]
  :=
  \left(
    (V_n)_{n=0}^{N+1},
    (Q_n)_{n=0}^{N};
    \mathfrak R_{N,M}(u;T)
  \right),
\]
with
\[
  V_n=\partial_t^nu,
  \qquad
  Q_n=\partial_t^np.
\]
Each retained temporal row obeys
\[
  V_{n+1}
  +
  \sum_{a=0}^{n}
  \binom{n}{a}
  (V_a\cdot\nabla)V_{n-a}
  +\nabla Q_n
  =
  \nu\Delta V_n,
\]
\[
  \nabla\cdot V_n
  =0,
\]
and
\[
  -\Delta Q_n
  =
  \partial_i\partial_j
  \sum_{a=0}^{n}
  \binom{n}{a}
  (V_a)_i(V_{n-a})_j.
\]
Rapid decay fixes every pressure constant.  At the retained depths, the mixed
velocity coordinates are the weak derivatives
\[
  J_{n,\alpha}=\partial_x^\alpha V_n,
\]
and the pressure coordinates are the corresponding derivatives of \(Q_n\).
Thus \(\mathfrak R_{N,M}\) measures the finite velocity block, whereas
\(\cR_{N,M}\) retains its pressure-complete equation identity.

To restrict such a block in time, let \(J=(0,\tau)\),
\(0<\tau\leq T_*\), and let \(\mathscr X_{N,M}(J)\) consist of records with
\[
  W_n\in L^2(J;H^{M+1}),
  \qquad
  W_{n+1}\in L^2(J;H^{M-1})
  \qquad(0\leq n\leq N),
\]
with \(\partial_tW_n=W_{n+1}\) distributionally on \(J\times\R^3\), and
normalized pressure rows
\[
  P_n
  =R_iR_j\sum_{a=0}^{n}\binom{n}{a}(W_a)_i(W_{n-a})_j
  \in L^1(J;H^M)
  \qquad(0\leq n\leq N),
\]
Their gradients lie in \(L^1(J;H^{M-1})\), and the displayed momentum and
divergence equations hold distributionally.  The scalar component is the
adjacent-row norm of the retained velocities on \(J\).  In particular,
\[
  \cR_{N,M}[u_0;T]
  \in
  \mathscr X_{N,M}((0,T)).
\]

As \((N,M)\) ranges over finite pairs, the rectangles cover every finite
coordinate of the intersection, each with its own finite bound.  Overlapping
rectangles contain the same distribution wherever they name the same
derivative.  That literal agreement is the datum needed for terminal assembly.
